\section{Type Ia supernovae}\label{sec:sne}

\subsection{Data}
We use the Pantheon+ sample \citep{Brout2022,Scolnic2022}: 1701 light curves of 1550 distinct SNe~Ia, of which 1590 light curves remain after the standard $z>0.01$ cut on $z_{\rm HD}$ (which removes the Cepheid-calibrated hosts and the peculiar-velocity-dominated tail), spanning $z=0.01$--$2.26$.
The distance modulus is $\mu=5\log_{10}D_L+25$ ($D_L$ in Mpc) with $D_L=(1+z_{\rm HEL})\,D_M(z_{\rm HD})$, the Pantheon+ standard cosmological convention in which the comoving-distance integral runs over the peculiar-velocity-corrected $z_{\rm HD}$ and only the $(1+z)$ factor uses the heliocentric redshift. The absolute magnitude $M$ is a free nuisance parameter, profiled jointly with $\Om$ for the \LCDM{} fit; for the bimetric model $\Om$ is implied by the adopted anchor through the quartic closure and held fixed.
We report two covariance treatments: the diagonal approximation (\code{m\_b\_corr\_err\_DIAG}), which neglects off-diagonal systematics, and the full STAT+SYS covariance ($1590\times1590$), which is primary.

As a cross-check we evaluate the non-standard variant with $z_{\rm CMB}$ in the distance integral (and the event selection also on $z_{\rm CMB}$, $N=1588$). It returns $\Dchi=6.50$ and moves the \LCDM{} best-fit $\Om$ from $0.331$ to $0.347$. Changing the integral alone (the $z_{\rm HD}$ selection retained) gives $6.33$ at $\Om=0.346$, while changing the selection alone gives $3.10$: the redshift entering the integral, not the event selection, carries the dominant share of the convention sensitivity.

\subsection{Results}\label{sec:sneres}
\paragraph{Diagonal covariance.} The \LCDM{} best fit is $\Om=0.350$ ($\chi^2/\mathrm{dof}=0.439$). The bimetric model ($\Om=0.30435$ implied by the adopted anchor) gives $\Dchi=7.32$ and a pure correction $\Dchi=1.01$ (evaluated at the bimetric $\Om$): the $E^2(z)$ correction alone is sub-detectable, and the sensitivity is carried by the $\Om$ mismatch.

\paragraph{Full covariance (STAT+SYS, primary).} The \LCDM{} best fit moves to $\Om=0.331$ ($\chi^2_{\min}=1402.92$, $\chi^2/\mathrm{dof}=0.883$); the bimetric $\chi^2=1405.86$ gives $\Dchi=2.94$ (standard) and $\Dchi_{\rm pure}=0.64$, with the mismatch component alone contributing $2.30$. The systematic correlations pull the \LCDM{} best-fit $\Om$ from $0.350$ towards the anchor-implied bimetric value, shrinking the mismatch that dominates $\Dchi$; the qualitative conclusion (SNe sensitive through the anchor-implied $\Om$) is unchanged.
The result is robust to the low-redshift cut. With the low-redshift cut $z>0.023$ of H\&M applied to the $z_{\rm CMB}$ column (1367 light curves) we obtain $\Dchi=2.69$ (applied to the $z_{\rm HD}$ column: 1371 light curves and $\Dchi=2.17$); selecting events by $z_{\rm CMB}$ instead of $z_{\rm HD}$ at $z>0.01$ moves $\Dchi$ from $2.94$ to $3.10$. No additional intrinsic dispersion is added to either covariance (the published errors already include it; the diagonal approximation neglects the off-diagonal systematics and here happens to yield a larger $\Dchi$).

The penalty is thus the known density split between the SNe-preferred $\Om$ and the rigid value implied by the adopted anchor through quartic closure, rather than a direct late-time shape mismatch (see \cref{sec:Om-H0}).

The $(\alpha,\beta)$ scan gives $\Dchi=7.315$ (diagonal) and $2.936$ (full) at every one of the 19 non-degenerate combinations, and the $\mFP/H_0$ scan ($15$ to $10^5$) is flat at the same values: the SNe sensitivity, like the anchor-implied $\Om$, is blind to the potential parameters across the grid. A $\theta$ scan at fixed $\OL=0.699$ with the self-consistent $\Om(\theta)$ gives $7.315$--$7.320$ (diagonal) and $2.936$--$2.939$ (full) over $6.1$--$31.5\deg$: over this scan the response is governed predominantly by the effective density $\cos^2\!\theta\,\Om\simeq1-\OL$, together with the tiny effective-law residual (\cref{sec:bao}). Square roots of $\Dchi$ are not calibrated as standalone $\sigma$ preferences and are not reported as such.
